\chapter{Streaming and Massive Graphs}
\label{ch:streaming}

Until now we have assumed the whole graph sits in memory, and that we may
poke at any vertex or edge whenever we like. That assumption breaks the
moment the graph is a web crawl with $10^{11}$ edges, a social network that
arrives as a firehose of friend-events, or a router log that streams past
faster than any disk can absorb. In those regimes the graph never fully
exists in RAM at one instant. We are handed its edges one at a time, in an
order we do not control, and we are allowed only a sliver of working memory
--- far less than the graph itself.

This chapter develops the \emph{streaming} model for graphs: what it costs,
what it buys, and (the part people forget) what it makes flatly impossible.
The headline regime is \emph{semi-streaming}: $\Oh(n\,\poly\log n)$ space,
enough to remember something about every vertex but nowhere near enough to
remember every edge. We will see that connectivity, bipartiteness, and
sparse spanners survive in that budget; that exact shortest paths and exact
triangle counts do not; and that a beautiful piece of linear-algebraic
trickery --- \emph{linear sketches} and $L_0$ samplers --- lets us handle
streams where edges are not only inserted but deleted.

\chapterline

\section{The streaming model, precisely}
\label{ch40:model}

\begin{definition}[Graph stream]\label{ch40:def-stream}
Let $\Gph=(V,E)$ have $n=\abs{V}$ vertices, known in advance, labelled
$1,\dots,n$. A \emph{graph stream} is a sequence
$\sigma = \langle a_1, a_2, \dots, a_m\rangle$ of \emph{tokens} presented one
at a time. In the \emph{insertion-only} (or \emph{cash-register}) model each
token $a_i$ is an edge $\set{u,v}$ to be added. In the \emph{dynamic} (or
\emph{turnstile}) model each token is a pair $(\set{u,v}, \pm 1)$ that
inserts or deletes an edge. An algorithm reads tokens left to right, keeps a
memory state $S$ of bounded size, and updates $S$ per token. After the last
token it answers a query about $\Gph$.
\end{definition}

Three knobs define the cost:

\begin{itemize}
\item \textbf{Space} $s$: bits of working memory, the scarce resource. Reading
a token does not count against $s$; only what we choose to retain does.
\item \textbf{Passes} $p$: how many times we are allowed to scan $\sigma$ from
the start. The default and most demanding setting is $p=1$.
\item \textbf{Per-token (update) time} and \textbf{post-processing time}: usually
secondary, but we track them.
\end{itemize}

The input has $n$ vertices and up to $m$ edges, with $m$ potentially as large
as $\binom{n}{2}=\Th(n^2)$. The whole point is that $s \ll m$.

\begin{definition}[Semi-streaming]\label{ch40:def-semistream}
An algorithm is \emph{semi-streaming} if it uses $\Oh(n\,\poly\log n)$ bits of
space and a small number of passes (ideally one).
\end{definition}

\begin{intuitionbox}
Why $n\,\poly\log n$ and not, say, $\sqrt{n}$ or $n^{1.5}$? Because $n\log n$
bits is exactly enough to store $\Oh(n)$ machine words --- one or a few per
vertex. That is the natural ``remember a little about everyone, forget the
edges'' frontier. Below it you cannot even hold a spanning forest; above it
($n^{1+c}$) you can hoard subquadratically many edges and the model stops
being interesting. Semi-streaming is the sweet spot where vertices fit but
edges do not, and almost every positive result in this chapter lives there.
\end{intuitionbox}

\begin{historybox}
The streaming model for numerical data was crystallized by Alon, Matias, and
Szegedy in 1996 (frequency moments, the AMS sketch). The graph-specific
\emph{semi-streaming} model is due to Feigenbaum, Kannan, McGregor, Suri, and
Zhang (2005), who pinned $\Oh(n\,\poly\log n)$ as the right budget and showed
connectivity and spanners fit inside it. The turnstile breakthrough ---
handling edge deletions via linear sketches --- came from Ahn, Guha, and
McGregor (2012), whose graph-sketching framework reduced connectivity to
$L_0$ sampling, recasting a combinatorial problem as pure linear algebra.
\end{historybox}

\subsection{A running graph to keep in mind}

Throughout, return to this little eight-vertex graph. Its edges will arrive
in a scrambled order so we can watch an online algorithm cope.

\begin{figure}[H]
\centering
\begin{tikzpicture}[
  every node/.style={circle,draw=nodeblue,fill=blue!6,minimum size=7mm,
    font=\small,inner sep=0pt},
  edge/.style={draw=edgegray,thick}]
\node (1) at (0,2) {1};
\node (2) at (2,2.6) {2};
\node (3) at (4,2) {3};
\node (4) at (1,0.4) {4};
\node (5) at (3,0.4) {5};
\node (6) at (5.2,1) {6};
\node (7) at (6.8,2) {7};
\node (8) at (6.8,0) {8};
\draw[edge] (1)--(2); \draw[edge] (2)--(3); \draw[edge] (1)--(4);
\draw[edge] (4)--(5); \draw[edge] (2)--(4); \draw[edge] (3)--(5);
\draw[edge] (3)--(6); \draw[edge] (6)--(7); \draw[edge] (6)--(8);
\draw[edge] (7)--(8);
\end{tikzpicture}
\caption{The running graph $\Gph$: eight vertices, ten edges. Vertices
$\set{1,2,3,4,5}$ form one cluster and $\set{6,7,8}$ a triangle, joined by the
single bridge $\set{3,6}$. There is no edge $3$--$4$, so the only triangle on
the left is $\set{1,2,4}$; the only triangle on the right is $\set{6,7,8}$.}
\label{ch40:fig-running}
\end{figure}

We will feed its edges in the adversarial order
\[
\sigma = \langle\,(3,6),(1,2),(6,7),(2,4),(4,5),(1,4),(2,3),(3,5),(7,8),(6,8)\,\rangle.
\]

\section{Connectivity in one pass}
\label{ch40:connectivity}

The cleanest semi-streaming result: decide whether $\Gph$ is connected (and
in fact maintain its connected components) in one pass and $\Oh(n\log n)$
bits.

\begin{intuitionbox}
A spanning forest already tells you the components, and a forest on $n$
vertices has at most $n-1$ edges. So we never need to store more than $n-1$
edges. Maintain a forest incrementally: when an edge $\set{u,v}$ arrives, add
it only if $u$ and $v$ are not already connected by what we have kept ---
otherwise it would close a cycle and is useless for connectivity. That is
exactly Kruskal's forest-building, run online, with a union--find structure
deciding ``already connected?'' in near-constant amortized time.
\end{intuitionbox}

\begin{algorithm}[H]
\caption{\textsc{StreamingSpanningForest}}
\label{ch40:alg-forest}
\DontPrintSemicolon
\KwIn{vertex set $V=\set{1,\dots,n}$; insertion stream of edges}
\KwOut{a spanning forest $F$; component labels}
initialize union--find $\mathcal{D}$ with each $v$ in its own set\;
$F \gets \emptyset$\;
\ForEach{arriving edge $\set{u,v}$}{
  \If{$\textsc{Find}(u)\neq\textsc{Find}(v)$}{
    $F \gets F \cup \set{\set{u,v}}$\;
    $\textsc{Union}(u,v)$\;
  }
}
\Return $F$ and the \textsc{Find} labels\;
\end{algorithm}

\begin{whybox}
\textbf{Correctness.} Claim: at every moment $F$ is a spanning forest of the
graph seen so far, with the \emph{same} connected components. Induction on
tokens. Initially both are empty/trivial. When $\set{u,v}$ arrives: if $u,v$
are already in the same tree of $F$, then by the induction hypothesis they
were already connected in the seen graph, so adding the edge changes no
component and skipping it changes no component --- $F$ stays a forest. If they
are in different trees, the seen graph just merged two components, and adding
$\set{u,v}$ to $F$ merges exactly those two trees without creating a cycle (it
joins two distinct trees). So $F$ remains a forest realizing the same
partition. After all tokens, $\Gph$ is connected iff $F$ is a single tree iff
$\abs{F}=n-1$.
\end{whybox}

\begin{complexitybox}
$F$ holds at most $n-1$ edges $=\Oh(n\log n)$ bits; union--find adds another
$\Oh(n\log n)$. Total space $\Oh(n\log n)$ --- semi-streaming with one pass.
Per-token work is one \textsc{Find}/\textsc{Union} pair, $\Oh(\alpha(n))$
amortized with path compression and union by rank. Post-processing is
$\Oh(n)$.
\end{complexitybox}

\begin{labbox}
Run Algorithm~\ref{ch40:alg-forest} on $\sigma$ from
Figure~\ref{ch40:fig-running}. Write the forest edges kept and the running
component count (start: 8 singletons).

\begin{center}
\begin{tabular}{llcc}
\toprule
token & action & kept? & \#comp \\
\midrule
$(3,6)$ & $3,6$ separate & yes & 7 \\
$(1,2)$ & separate & yes & 6 \\
$(6,7)$ & separate & yes & 5 \\
$(2,4)$ & separate & yes & 4 \\
$(4,5)$ & separate & yes & 3 \\
$(1,4)$ & $1,4$ already linked ($1$--$2$--$4$) & \textbf{no} & 3 \\
$(2,3)$ & joins $\set{1,2,4,5}$ with $\set{3,6,7}$ & yes & 2 \\
$(3,5)$ & $3,5$ already linked & \textbf{no} & 2 \\
$(7,8)$ & joins $8$ in & yes & 1 \\
$(6,8)$ & already linked & \textbf{no} & 1 \\
\bottomrule
\end{tabular}
\end{center}

We kept exactly $7=n-1$ edges and ended at one component: $\Gph$ is connected.
The three edges we dropped, $(1,4),(3,5),(6,8)$, are precisely the cycle-closing
ones --- the ``extra'' density the forest is allowed to forget.
\end{labbox}

\subsection{Bipartiteness and $k$-edge-connectivity for free-ish}

The same skeleton answers more questions with a constant factor more space.

\textbf{Bipartiteness.} Keep the spanning forest \emph{plus} a parity bit
$\phi(v)\in\set{0,1}$ per vertex (the side of the 2-coloring relative to its
tree root, maintained as union--find merges). When a non-tree edge
$\set{u,v}$ arrives with $\textsc{Find}(u)=\textsc{Find}(v)$, check whether
$\phi(u)=\phi(v)$; if so we have found an odd cycle and the graph is not
bipartite. This is one pass, $\Oh(n\log n)$ bits.

\textbf{$k$-edge-connectivity.} Maintain $k$ edge-disjoint forests
$F_1,\dots,F_k$ (Nagamochi--Ibaraki / Thurimella sparsification): insert each
arriving edge into the first $F_i$ where it does not close a cycle. The union
$F_1\cup\cdots\cup F_k$ is a \emph{$k$-skeleton} preserving all cuts of size
$\le k$, so it answers $k$-edge-connectivity exactly. Space $\Oh(kn\log n)$:
semi-streaming for $k=\poly\log n$.

\begin{connectionbox}
The forest trick is Kruskal (Chapter~6) without sorting: the stream supplies
an arbitrary order and we keep whatever extends the forest. The $k$-forest
sparsifier is the combinatorial cousin of the \emph{spectral} sparsifier from
Chapters~31--34: there we preserved every quadratic form $x^\top\Lap x$ to
within $1\pm\eps$ using $\Oh(n\log n/\eps^2)$ \emph{weighted} edges; here we
preserve every cut of value $\le k$ exactly using $kn$ \emph{unweighted}
edges. Both say the same moral: a graph's connectivity structure is
compressible to near-linear size.
\end{connectionbox}

\section{Spanners: approximate distances in small space}
\label{ch40:spanners}

A spanning forest forgets distances completely --- it can stretch a length-1
edge into a length-$\Th(n)$ detour. If we want approximate shortest paths in
one pass, we need a \emph{spanner}.

\begin{definition}[Spanner]\label{ch40:def-spanner}
A subgraph $H\subseteq\Gph$ is a \emph{$t$-spanner} (stretch $t$) if for all
$u,v$,
\[
\dist_\Gph(u,v) \;\le\; \dist_H(u,v) \;\le\; t\cdot \dist_\Gph(u,v).
\]
Equivalently (it suffices to check edges): for every edge $\set{u,v}\in\Gph$,
$\dist_H(u,v)\le t$.
\end{definition}

The ``edges suffice'' remark is worth a line: if every edge is $t$-stretched,
then concatenating along any shortest path stretches the whole path by at most
$t$. So a streaming spanner only has to guarantee, edge by edge, that it never
drops an edge whose endpoints are currently far apart in $H$.

\begin{intuitionbox}
Greedy spanner, online: when $\set{u,v}$ arrives, ask ``are $u$ and $v$
already within distance $t$ using the edges I have kept?'' If yes, the new
edge is redundant --- drop it. If no, keep it. This guarantees stretch $t$ by
construction. The payoff is that the kept graph $H$ then has \emph{high girth}:
every cycle in $H$ has length $> t+1$, because the last edge of any short
cycle would have been dropped. High girth forces sparsity.
\end{intuitionbox}

\begin{algorithm}[H]
\caption{\textsc{StreamingGreedySpanner}$(t)$, $t=2k-1$ odd}
\label{ch40:alg-spanner}
\DontPrintSemicolon
\KwIn{stream of edges; stretch parameter $t=2k-1$}
\KwOut{a $t$-spanner $H$}
$H \gets \emptyset$\;
\ForEach{arriving edge $\set{u,v}$}{
  compute $d \gets \dist_H(u,v)$ by BFS in $H$ (stop early at depth $t$)\;
  \If{$d > t$ \textbf{(includes $d=\infty$)}}{
    $H \gets H \cup \set{\set{u,v}}$\;
  }
}
\Return $H$\;
\end{algorithm}

\begin{whybox}
\textbf{Stretch.} For any kept edge, directly $\dist_H\le 1\le t$. For any
\emph{dropped} edge $\set{u,v}$, we dropped it precisely because
$\dist_H(u,v)\le t$ at that moment, and $H$ only grows, so the bound persists.
Hence every original edge is $t$-stretched, and by the edge-suffices remark
$H$ is a $t$-spanner.

\textbf{Sparsity via girth.} Suppose $H$ contained a cycle $C$ of length
$\ell\le t+1=2k$. Consider the \emph{last} edge $\set{u,v}$ of $C$ to be added
to $H$. At that instant the other $\ell-1\le t$ edges of $C$ were already in
$H$, forming a $u$--$v$ path of length $\ell-1\le t$, so $\dist_H(u,v)\le t$
and we would have \emph{dropped} $\set{u,v}$ --- contradiction. Therefore $H$
has girth $\ge t+2 = 2k+1$.
\end{whybox}

\begin{lemma}[High-girth graphs are sparse, the Moore bound]\label{ch40:lem-moore}
A graph on $n$ vertices with girth $\ge 2k+1$ has $\Oh(n^{1+1/k})$ edges.
\end{lemma}

\begin{proof}[Proof sketch]
If the average degree is $d=2\abs{E}/n$, a standard cleaning gives a subgraph
of minimum degree $\ge d/2$. Pick any vertex and grow a BFS tree to depth
$k$; girth $\ge 2k+1$ means no two of these short paths collide before depth
$k$, so the ball is a genuine tree and contains at least
$1+(d/2)+(d/2)^2+\cdots$ down to depth $k$, i.e.\ $\gtrsim (d/2)^k$ vertices.
This must be $\le n$, giving $d/2 \le n^{1/k}$, so
$\abs{E}=nd/2 \le n\cdot n^{1/k}=n^{1+1/k}$ up to the additive cleaning term.
\end{proof}

\begin{complexitybox}
Algorithm~\ref{ch40:alg-spanner} with $t=2k-1$ stores an
$\Oh(n^{1+1/k})$-edge $(2k-1)$-spanner. Choosing $k=\log n$ gives stretch
$\Oh(\log n)$ and space $\Oh(n^{1+1/\log n}\log n)=\Oh(n\log n)$ bits ---
semi-streaming. The per-token cost is one truncated BFS, $\Oh(\abs{E(H)})$ in
the worst case; the additive-error and faster variants of Baswana--Sen reduce
this to $\Oh(1)$ expected per edge.
\end{complexitybox}

\begin{labbox}
Build a $3$-spanner ($t=3$, $k=2$) of the running graph in stream order
$\sigma$. Rule: keep $\set{u,v}$ unless $\dist_H(u,v)\le 3$ already.

\begin{center}
\begin{tabular}{llc}
\toprule
token & current $\dist_H(u,v)$ & keep? \\
\midrule
$(3,6)$ & $\infty$ & yes \\
$(1,2)$ & $\infty$ & yes \\
$(6,7)$ & $\infty$ & yes \\
$(2,4)$ & $\infty$ & yes \\
$(4,5)$ & $\infty$ & yes \\
$(1,4)$ & $1$--$2$--$4$ $=2\le3$ & \textbf{no} \\
$(2,3)$ & $\infty$ ($\set{1,2,4,5}$ vs $\set{3,6,7}$) & yes \\
$(3,5)$ & $3$--$2$--$4$--$5$ $=3\le3$ & \textbf{no} \\
$(7,8)$ & $\infty$ & yes \\
$(6,8)$ & $6$--$7$--$8$ $=2\le3$ & \textbf{no} \\
\bottomrule
\end{tabular}
\end{center}

Kept: $\set{(3,6),(1,2),(6,7),(2,4),(4,5),(2,3),(7,8)}$ --- seven edges, a
tree here (the graph is small so the spanner coincides with a forest). On a
denser instance the spanner would keep cycle edges that a forest discards, in
exchange for the stretch guarantee. Sanity check on the dropped edge $(3,5)$:
in $H$, $\dist(3,5)=3=3\cdot\dist_\Gph(3,5)$, exactly the stretch ceiling.
Figure~\ref{ch40:fig-spanner} draws the kept and dropped edges.
\end{labbox}

\begin{figure}[H]
\centering
\begin{tikzpicture}[
  scale=0.9,
  token/.style={rectangle,draw=gray!40,fill=white,rounded corners=2pt,font=\tiny\ttfamily,inner sep=2pt},
  vnode/.style={circle,draw=nodeblue,fill=blue!6,minimum size=6mm,font=\scriptsize,inner sep=0pt},
  kept/.style={draw=targetgreen,very thick},
  dropped/.style={draw=trapred,thick,dashed}
]
  % Lane 1: Stream Timeline
  \fill[sourceblue!5,rounded corners=3pt,draw=sourceblue!25] (-0.5, 4.0) rectangle (12.5, 5.8);
  \node[font=\footnotesize\scshape,text=sourceblue!90!black,anchor=west] at (-0.3, 5.5) {Stream Timeline $\sigma$: Online Edge Decisions};
  
  \node[token] (t1) at (0.2, 4.6) {(3,6)};
  \node[above=1pt of t1,font=\tiny,text=targetgreen!80!black] {keep};
  
  \node[token] (t2) at (1.4, 4.6) {(1,2)};
  \node[above=1pt of t2,font=\tiny,text=targetgreen!80!black] {keep};
  
  \node[token] (t3) at (2.6, 4.6) {(6,7)};
  \node[above=1pt of t3,font=\tiny,text=targetgreen!80!black] {keep};
  
  \node[token] (t4) at (3.8, 4.6) {(2,4)};
  \node[above=1pt of t4,font=\tiny,text=targetgreen!80!black] {keep};
  
  \node[token] (t5) at (5.0, 4.6) {(4,5)};
  \node[above=1pt of t5,font=\tiny,text=targetgreen!80!black] {keep};
  
  \node[token,draw=trapred!60,fill=trapred!5] (t6) at (6.2, 4.6) {(1,4)};
  \node[above=1pt of t6,font=\tiny\bfseries,text=trapred] {drop ($d{=}2$)};
  
  \node[token] (t7) at (7.4, 4.6) {(2,3)};
  \node[above=1pt of t7,font=\tiny,text=targetgreen!80!black] {keep};
  
  \node[token,draw=trapred!60,fill=trapred!5] (t8) at (8.6, 4.6) {(3,5)};
  \node[above=1pt of t8,font=\tiny\bfseries,text=trapred] {drop ($d{=}3$)};
  
  \node[token] (t9) at (9.8, 4.6) {(7,8)};
  \node[above=1pt of t9,font=\tiny,text=targetgreen!80!black] {keep};
  
  \node[token,draw=trapred!60,fill=trapred!5] (t10) at (11.2, 4.6) {(6,8)};
  \node[above=1pt of t10,font=\tiny\bfseries,text=trapred] {drop ($d{=}2$)};

  \draw[->,gray!60,thick] (-0.3, 4.2) -- (12.3, 4.2) node[right,font=\tiny]{time $t$};

  % Lane 2: Spanner Topology
  \fill[gray!4,rounded corners=3pt,draw=gray!25] (-0.5, 0.2) rectangle (12.5, 3.6);
  \node[font=\footnotesize\scshape,text=gray!80!black,anchor=west] at (-0.3, 3.3) {Maintained Spanner Graph $H \subseteq \Gph$ (Stretch $2k-1 = 3$)};

  \begin{scope}[shift={(1.5, 0.4)}]
    \node[vnode] (1) at (0,1.8) {1};
    \node[vnode] (2) at (1.8,2.4) {2};
    \node[vnode] (3) at (3.6,1.8) {3};
    \node[vnode] (4) at (0.9,0.5) {4};
    \node[vnode] (5) at (2.7,0.5) {5};
    \node[vnode] (6) at (5.0,1.2) {6};
    \node[vnode] (7) at (6.6,2.0) {7};
    \node[vnode] (8) at (6.6,0.4) {8};

    \draw[kept] (3)--(6); \draw[kept] (1)--(2); \draw[kept] (6)--(7);
    \draw[kept] (2)--(4); \draw[kept] (4)--(5); \draw[kept] (2)--(3);
    \draw[kept] (7)--(8);

    \draw[dropped] (1)--(4); \draw[dropped] (3)--(5); \draw[dropped] (6)--(8);
  \end{scope}

  \begin{scope}[shift={(9.6, 1.8)}]
    \draw[kept] (0,0.5) -- (0.8,0.5) node[right,font=\tiny]{kept ($7$ edges)};
    \draw[dropped] (0,0.0) -- (0.8,0.0) node[right,font=\tiny]{dropped ($3$ edges)};
    \node[font=\tiny,align=left,text=gray!70!black] at (0.8,-0.7) {dropped because\\$\dist_H(u,v) \le 3$};
  \end{scope}
\end{tikzpicture}
\caption{Semi-streaming greedy construction of a $3$-spanner on the running graph under adversarial stream order $\sigma$. Top lane: sequential stream arrivals. Each token $\{u,v\}$ is inspected; if $\dist_H(u,v) \le 3$ via existing edges, it is dropped as redundant, otherwise it is retained. Bottom lane: the resulting spanner $H$. Solid green edges form $H$ ($7$ edges, preserving connectivity and bounding stretch to $\le 3$); dashed red edges $\set{1,4},\set{3,5},\set{6,8}$ were pruned with zero memory penalty.}
\label{ch40:fig-spanner}
\end{figure}

\begin{trapbox}
``A spanning forest is a spanner.'' It is a $1$-spanner of nothing useful: a
forest preserves \emph{connectivity} but can blow up \emph{distance}
arbitrarily. Connectivity and distance are different currencies; do not pay
for one expecting the other.
\end{trapbox}

\section{Dynamic streams: when edges get deleted}
\label{ch40:sketching}

Algorithms~\ref{ch40:alg-forest} and~\ref{ch40:alg-spanner} are hopeless
under deletions. The spanning forest may keep edge $\set{u,v}$, and later a
token $\bigl(\set{u,v},-1\bigr)$ removes it --- but the forest has no idea
which \emph{other} edge should take its place, because it deliberately threw
those edges away. Greedy retention is fundamentally non-reversible.

The fix is to store, instead of a chosen subset of edges, a small
\emph{linear} summary of the entire edge multiset. Linearity is exactly what
makes deletions painless: a deletion is just an insertion of $-1$, and the
summary updates by subtraction.

\subsection{The vertex--edge incidence vector}

Index the coordinate space by unordered pairs $\binom{V}{2}$, of dimension
$\binom{n}{2}$. For a vertex $v$, define its \emph{signed incidence vector}
$x_v \in \R^{\binom{n}{2}}$: for each edge $\set{u,v}$ present in the graph,
$x_v$ has a $+1$ or $-1$ in coordinate $\set{u,v}$, with the sign chosen by a
fixed global rule, say $+1$ if $v<u$ and $-1$ if $v>u$.

\begin{lemma}[Boundary cancellation]\label{ch40:lem-boundary}
For any vertex subset $S\subseteq V$, the support of $\sum_{v\in S} x_v$ is
exactly the set of edges crossing the cut $(S,\bar S)$.
\end{lemma}

\begin{proof}
Coordinate $\set{u,w}$ receives a contribution from $x_u$ and from $x_w$ only.
If both $u,w\in S$, the two contributions have opposite signs (one endpoint
sees the edge as ``smaller'', the other as ``larger'') and cancel to $0$. If
both are outside $S$, neither contributes. If exactly one is in $S$, the
coordinate survives with value $\pm1$. So the surviving coordinates are
precisely the cut edges of $S$.
\end{proof}

\begin{intuitionbox}
This is the incidence-matrix identity $\Lap = B^\top B$ from
Chapter~31 wearing streaming clothes. Summing rows of the incidence matrix
over $S$ cancels internal edges and leaves the boundary $\partial S$. So if we
could, for any $S$, recover \emph{one} nonzero coordinate of $\sum_{v\in
S}x_v$, we would learn one edge leaving $S$ --- enough to grow a spanning
forest by repeatedly merging components and asking for an outgoing edge. The
catch: $\sum_{v\in S}x_v$ lives in $\binom{n}{2}$ dimensions and we cannot
store it. We store a tiny \emph{sketch} of each $x_v$ instead, and exploit
that the sum of sketches is the sketch of the sum.
\end{intuitionbox}

\subsection{$L_0$ sampling}

The primitive we need is a sampler that, from a high-dimensional vector
accessed only through a small linear sketch, returns a uniformly random
nonzero coordinate.

\begin{definition}[$L_0$ sampler]\label{ch40:def-l0}
An $L_0$ sampler for vectors $x\in\R^N$ is a randomized linear map
$\mathsf{Sk}:\R^N\to\R^s$ together with a recovery procedure $\mathsf{Rec}$
such that, from $\mathsf{Sk}(x)$ alone, with probability $\ge 1-\delta$,
$\mathsf{Rec}$ outputs a coordinate $i$ with $x_i\neq 0$, and conditioned on
success $i$ is (near-)uniform over $\mathrm{supp}(x)$. Such samplers exist with
$s=\Oh(\log^2 N \cdot \log(1/\delta))$.
\end{definition}

The standard construction \emph{subsamples by levels}. Hash each coordinate to
a geometric level $j$ with $\Pr[\text{level}\ge j]=2^{-j}$. At level
$j\approx \log\norm{x}_0$ the survivors number $\Th(1)$ in expectation; at
that level keep a tiny linear ``$1$-sparse recovery'' gadget that can decode a
single nonzero coordinate exactly. Run all $\Oh(\log N)$ levels in parallel,
read off the first level that is $1$-sparse. Each gadget is linear, so the
whole sampler is linear, so deletions are handled by subtraction. A
$1$-sparse recovery gadget stores three running linear combinations
$\bigl(\sum_i x_i,\ \sum_i i\,x_i,\ \sum_i \mathrm{fp}(i)\,x_i\bigr)$;
when the support is a single index $i^\star$ with value $a$, the first slot is
$a$, the second is $a\,i^\star$, and dividing recovers $i^\star=$ (slot2)/(slot1),
with the fingerprint slot detecting (whp) that the support really was size one.

\begin{whybox}
\textbf{Why it is exactly linear, and why that matters.} Every quantity stored
--- the level hashes are fixed in advance, and each gadget keeps only
\emph{sums} $\sum_i c_i x_i$ for fixed coefficients $c_i$. When token
$(\set{u,v},\pm1)$ arrives we add $\pm1$ to two incidence vectors $x_u,x_v$,
which adds $\pm c_{\set{u,v}}$ into the relevant sums. Insertion and deletion
are the same operation with opposite sign. The sketch never ``remembers
which edges were kept''; it remembers numerical aggregates that automatically
forget a deleted edge. This is the entire reason linear sketches survive
turnstile streams while greedy forests cannot.
\end{whybox}

\subsection{Connectivity in dynamic streams (AGM sketching)}

Now assemble the pieces. Maintain, for each vertex $v$, an independent
collection of $\Oh(\log n)$ $L_0$ sketches of its incidence vector $x_v$ ---
one fresh sketch per round of a Bor\r{u}vka-style component merge.

\begin{algorithm}[H]
\caption{\textsc{SketchConnectivity}}
\label{ch40:alg-agm}
\DontPrintSemicolon
\KwIn{turnstile edge stream; $R=\ceil{\log_2 n}+1$ rounds}
\KwOut{a spanning forest of the final graph}
\textbf{During the stream:} for each $v$ keep $R$ independent $L_0$ sketches
$\mathsf{Sk}_1(x_v),\dots,\mathsf{Sk}_R(x_v)$, updated linearly per token\;
\textbf{After the stream (Bor\r{u}vka):}\;
$\mathcal{C}\gets$ partition of $V$ into singletons; $F\gets\emptyset$\;
\For{$r=1$ \KwTo $R$}{
  \ForEach{component $S\in\mathcal{C}$}{
    using \emph{round-$r$} sketches, compute
    $\mathsf{Sk}_r\!\bigl(\sum_{v\in S}x_v\bigr)=\sum_{v\in S}\mathsf{Sk}_r(x_v)$\;
    run $\mathsf{Rec}$ to sample an edge leaving $S$ (a nonzero coordinate)\;
  }
  add every sampled inter-component edge to $F$ and merge the components it
  joins, updating $\mathcal{C}$\;
}
\Return $F$\;
\end{algorithm}

\begin{whybox}
\textbf{Why $R=\Oh(\log n)$ rounds suffice.} Each Bor\r{u}vka round, every
current component samples (at least) one outgoing edge if one exists, so every
non-maximal component merges with another. The number of components therefore
at least halves each round; after $\ceil{\log_2 n}$ rounds it stops shrinking,
and the partition equals the true connected components. Each round must use a
\emph{fresh independent} sketch: the components queried in round $r$ depend on
the random outcomes of rounds $1,\dots,r-1$, so reusing a sketch would
correlate the query set with the sketch's randomness and void the $L_0$
guarantee. That is why we pay $\Oh(\log n)$ sketches per vertex, not one.

\textbf{Why the sum-of-sketches step is valid.} By linearity
$\sum_{v\in S}\mathsf{Sk}_r(x_v)=\mathsf{Sk}_r(\sum_{v\in S}x_v)$, and by
Lemma~\ref{ch40:lem-boundary} that argument vector is supported exactly on the
cut $\partial S$. An $L_0$ sample from it is therefore a (uniform) edge
leaving $S$ --- precisely Bor\r{u}vka's needed primitive.
\end{whybox}

\begin{figure}[H]
\centering
\begin{tikzpicture}[
  comp/.style={draw=structurepurple,thick,rounded corners,inner sep=3pt},
  every node/.style={circle,draw=nodeblue,fill=blue!6,minimum size=5.5mm,
    font=\scriptsize,inner sep=0pt},
  cut/.style={draw=trapred,very thick,dashed}]
% left: round start, components A={a,b}, B={c}, C={d,e}
\node (a) at (0,1.2) {a};
\node (b) at (0,0) {b};
\node (c) at (2.0,0.6) {c};
\node (d) at (4.0,1.2) {d};
\node (e) at (4.0,0) {e};
\draw[edgegray,thick] (a)--(b);
\draw[edgegray,thick] (d)--(e);
\draw[cut] (b)--(c);
\draw[cut] (c)--(d);
\begin{scope}[on background layer]
  \node[comp,fit=(a)(b),label=above:{\scriptsize $S$}] {};
  \node[comp,fit=(c)] {};
  \node[comp,fit=(d)(e)] {};
\end{scope}
\node[draw=none,fill=none,font=\scriptsize\itshape] at (2.0,-1.2)
  {$L_0$ sample from $\sum_{v\in S}x_v$ returns a dashed (cut) edge};
\end{tikzpicture}
\caption{One Bor\r{u}vka round under sketching. Summing the incidence sketches
over a component $S$ cancels its internal solid edges
(Lemma~\ref{ch40:lem-boundary}) and leaves only the dashed boundary edges; the
$L_0$ sampler returns one of them, merging $S$ with a neighbor.}
\label{ch40:fig-boruvka}
\end{figure}

\begin{complexitybox}
To make every one of the $\Oh(n\log n)$ component-queries across all rounds
succeed simultaneously, take the per-sketch failure to $\delta=1/\poly(n)$, so
each $L_0$ sketch is $\Oh(\log^2 N\cdot\log(1/\delta))=\Oh(\log^3 n)$ bits
(Definition~\ref{ch40:def-l0} with $N=\binom n2$, $\log N=\Th(\log n)$). Per
vertex we keep $R=\Oh(\log n)$ of them, giving $\Oh(\log^4 n)$ bits per vertex
and $\Oh(n\log^4 n)$ bits total --- semi-streaming, one pass, and it handles
arbitrary insertions and deletions. Sharper $L_0$ samplers shave this to the
$\Oh(n\log^3 n)$ headline figure of the original AGM analysis. The classical
insertion-only forest (Algorithm~\ref{ch40:alg-forest}) is cheaper at
$\Oh(n\log n)$; sketching pays the extra $\log$ factors precisely to buy
deletion-tolerance.
\end{complexitybox}

\begin{connectionbox}
Lemma~\ref{ch40:lem-boundary} is the streaming face of the boundary operator
$B$ (the signed incidence matrix) studied in the Laplacian chapter:
$\sum_{v\in S}x_v = B^\top \mathbf{1}_S$ is the indicator of the cut, and
$\norm{B^\top\mathbf{1}_S}_0=\abs{\partial S}$. Graph sketching is, quite
literally, doing Bor\r{u}vka's MST/connectivity algorithm
(Chapter~6) over a compressed coordinate space, with the cut/boundary algebra
of Chapter~31 supplying the correctness. The same incidence-sketch idea
extends (Ahn--Guha--McGregor) to spanners, sparsifiers, and even spectral
sparsifiers in turnstile streams.
\end{connectionbox}

\section{Approximate triangle counting}
\label{ch40:triangles}

Let $T$ be the number of triangles in $\Gph$. Triangle counts measure
clustering --- the tendency of friends-of-friends to be friends --- and feed
network-science statistics like the global clustering coefficient. Computing
$T$ exactly in small space is impossible (Section~\ref{ch40:lower}); we
estimate it.

\subsection{A clean sampling estimator (one pass)}

\begin{intuitionbox}
A triangle is three edges $\set{a,b},\set{b,c},\set{a,c}$. Pick one edge at
random as the ``base'', then a third vertex at random, and \emph{watch} the
rest of the stream to see whether the two edges that would complete the
triangle on that vertex actually show up. If we structure the random choice so
that each triangle has a known probability of being ``set up'' and then
``confirmed'', the indicator of confirmation is an unbiased $1/p$-scaled
estimator of $T$.
\end{intuitionbox}

A concrete one-pass estimator (Buriol et al.): sample a uniformly random
position $i$ in the stream and a uniformly random vertex $w\notin$ the edge
$e_i=\set{u,v}$ at that position. Then scan the \emph{rest} of the stream; let
the indicator $X=1$ if both edges $\set{u,w}$ and $\set{v,w}$ appear after
position $i$. Each triangle $\set{u,v,w}$ is detected exactly when we sampled
its earliest edge as $e_i$ and picked the opposite vertex $w$, an event of
probability $\tfrac{1}{m}\cdot\tfrac{1}{n-2}$. Hence
\[
\mathbb{E}[X] \;=\; \frac{T}{m\,(n-2)},
\qquad\text{so}\qquad
\widehat T \;=\; m\,(n-2)\,X
\]
is unbiased: $\mathbb{E}[\widehat T]=T$. Averaging $s$ independent copies
and applying Chebyshev drives the variance down.

\begin{algorithm}[H]
\caption{\textsc{TriangleEstimate} (one independent copy)}
\label{ch40:alg-tri}
\DontPrintSemicolon
\KwIn{insertion stream $e_1,\dots,e_m$; $n,m$}
\KwOut{unbiased estimate $\widehat T$ of triangle count}
pick uniform $i\in\set{1,\dots,m}$ via reservoir sampling\;
\tcp{when token $e_i=\set{u,v}$ is the current reservoir choice:}
pick uniform $w\in V\setminus\set{u,v}$; set $a\gets b\gets \text{false}$\;
\ForEach{token $e_j$ with $j>i$}{
  \lIf{$e_j=\set{u,w}$}{$a\gets\text{true}$}
  \lIf{$e_j=\set{v,w}$}{$b\gets\text{true}$}
}
$X\gets [\,a \wedge b\,]$\;
\Return $\widehat T \gets m\,(n-2)\,X$\;
\end{algorithm}

\begin{whybox}
\textbf{Unbiasedness.} Fix a triangle $\tau=\set{u,v,w}$ and let $e_i$ be its
first edge in stream order, say $\set{u,v}$. The copy ``catches'' $\tau$ iff
the reservoir landed on that very position \emph{and} the random third vertex
equals the remaining vertex $w$. These choices are uniform and independent, so
$\Pr[\text{catch }\tau]=\frac1m\cdot\frac1{n-2}$, and given a catch both
completing edges appear later (by definition of ``first edge''), so $X=1$.
Catches of distinct triangles are disjoint events for a single copy, hence
$\mathbb{E}[X]=\sum_\tau \frac1{m(n-2)}=\frac{T}{m(n-2)}$ and
$\mathbb{E}[\widehat T]=T$.
\end{whybox}

\begin{complexitybox}
One copy uses $\Oh(1)$ words (reservoir index, the chosen edge, $w$, two
bits). Reducing variance to get a $(1\pm\eps)$ estimate with confidence
$1-\delta$ requires
$s=\Oh\!\bigl(\tfrac{1}{\eps^2}\cdot\tfrac{m\,\Delta}{T}\cdot\log\tfrac1\delta\bigr)$
parallel copies, where $\Delta$ bounds how many triangles share an edge ---
intuitively, the harder the instance hides its triangles, the more samples you
must spend. With $T$ large relative to $m$ this is genuinely sublinear; with
$T$ tiny it degrades, which the lower bound below shows is unavoidable.
\end{complexitybox}

\begin{labbox}
The running graph has exactly two triangles: $\set{1,2,4}$ and $\set{6,7,8}$,
so $T=2$, with $m=10$, $n=8$, giving $\widehat T=10\cdot 6\cdot X=60X$ per
copy. The triangle $\set{6,7,8}$ has edges at positions $3\,(6,7)$, $9\,(7,8)$,
$10\,(6,8)$; its first edge is $(6,7)$ at position $3$. A copy catches it iff
the reservoir picks $i=3$ \emph{and} $w=8$: probability
$\frac1{10}\cdot\frac16=\frac1{60}$, contributing $60\cdot\frac1{60}=1$ to
$\mathbb{E}[\widehat T]$. The triangle $\set{1,2,4}$ contributes another $1$.
Sum: $\mathbb{E}[\widehat T]=2=T$. A single copy almost always reports $0$ and
occasionally $60$; only the \emph{average} of many copies is informative ---
the working lesson of high-variance estimators.
\end{labbox}

\begin{applicationbox}
Triangle and clustering-coefficient counts on streamed graphs drive spam and
bot detection (fraud rings form unusually dense triangle clusters),
recommendation (``triadic closure'' predicts new links), and bioinformatics
(motif counts in protein interaction networks). Real systems --- think a
social platform ingesting billions of edge events per day --- cannot store the
graph, so unbiased streaming estimators like
Algorithm~\ref{ch40:alg-tri} (and its sketch-based turnstile cousins) are how
these statistics get computed at all.
\end{applicationbox}

\subsection{Frequency moments and the AMS sketch}
\label{ch40:ams}

Beyond reservoir sampling on edges, algebraic linear sketching offers a unified framework for tracking
frequencies, degrees, and higher-order subgraph counts in streams with both insertions and deletions.
The foundational construction is the Alon--Matias--Szegedy (AMS) sketch ($1996$).

\begin{intuitionbox}
Suppose each edge $e$ in the universe has an accumulated frequency $f_e$ (in simple graphs $f_e \in \{0,1\}$,
while in dynamic multi-graphs $f_e \in \mathbb{Z}$). The second frequency moment $F_2 = \sum_e f_e^2$ measures
edge energy and degree variances.

To maintain $F_2$ in logarithmic space, choose a $4$-wise independent hash function $h: E \to \{-1, +1\}$
drawn from a family with $\mathbb{E}[h(e)] = 0$ and $\mathbb{E}[h(e)^2] = 1$. Maintain a single scalar
accumulator $Z$ initialized to $0$. Upon receiving an update $(e, \Delta)$ where $\Delta \in \{+1, -1\}$, update
\[
  Z \;\gets\; Z + \Delta \cdot h(e).
\]
At stream end, $Z = \sum_e f_e h(e)$. Its square $Y = Z^2$ is an unbiased estimator of $F_2$:
\[
  \mathbb{E}[Y] = \mathbb{E}\Bigl[\sum_e f_e^2 h(e)^2 + \sum_{e \neq e'} f_e f_{e'} h(e) h(e')\Bigr]
  = \sum_e f_e^2 \cdot 1 + 0 = F_2.
\]
By $4$-wise independence, pairwise cross-terms vanish and $\mathrm{Var}[Y] \le 2 F_2^2$.
To drive the estimation error down to $\pm \eps F_2$, we compute the average of $w = \lceil 2/\eps^2 \rceil$
independent accumulators. To drive the failure probability below $\delta$, we take the median of $d = \lceil 36\ln(2/\delta)\rceil$
such row averages. Because each cell is a linear combination of stream tokens, edge deletions are handled natively
with no memory overhead.
\end{intuitionbox}

\begin{figure}[H]
\centering
\begin{tikzpicture}[
  scale=0.9,
  token/.style={rectangle,draw=sourceblue!70!black,fill=sourceblue!10,rounded corners=2pt,font=\scriptsize\ttfamily,inner sep=3pt},
  hashnode/.style={rectangle,draw=quotientorange!80!black,fill=quotientorange!15,rounded corners=2pt,font=\scriptsize\bfseries,inner sep=3pt},
  reg/.style={rectangle,draw=targetgreen!80!black,fill=targetgreen!15,rounded corners=2pt,font=\scriptsize\bfseries,minimum width=1.1cm,minimum height=5mm},
  lanehead/.style={font=\footnotesize\scshape,anchor=west}
]
  % Lane 1: Stream Updates
  \fill[sourceblue!5,rounded corners=4pt,draw=sourceblue!25] (-0.8, 3.8) rectangle (12.2, 5.4);
  \node[lanehead,text=sourceblue!90!black] at (-0.6, 5.0) {Lane 1: Stream Updates $(e_t, \Delta_t \in \{\pm 1\})$};
  \node[token] (u1) at (1.0, 4.3) {$(e_1, +1)$};
  \node[token] (u2) at (3.5, 4.3) {$(e_2, +1)$};
  \node[token] (u3) at (6.0, 4.3) {$(e_1, -1)$};
  \node[token] (u4) at (8.5, 4.3) {$(e_3, +1)$};
  \draw[->,sourceblue!80!black,thick] (0, 4.0) -- (11.5, 4.0) node[right,font=\tiny]{stream order};

  % Lane 2: 4-wise Independent Hash Functions
  \fill[quotientorange!5,rounded corners=4pt,draw=quotientorange!25] (-0.8, 1.8) rectangle (12.2, 3.4);
  \node[lanehead,text=quotientorange!90!black] at (-0.6, 3.0) {Lane 2: $4$-Wise Independent Signs $h_r: E \to \{-1, +1\}$};
  \node[hashnode] (h1) at (1.5, 2.3) {Hash $h_1(e) \in \pm 1$};
  \node[hashnode] (h2) at (5.5, 2.3) {Hash $h_2(e) \in \pm 1$};
  \node[hashnode] (hd) at (9.5, 2.3) {Hash $h_d(e) \in \pm 1$};

  \draw[->,gray!60,thick] (3.5, 4.0) -- (1.5, 2.7);
  \draw[->,gray!60,thick] (3.5, 4.0) -- (5.5, 2.7);
  \draw[->,gray!60,thick] (3.5, 4.0) -- (9.5, 2.7);

  % Lane 3: Parallel Accumulators & Median-of-Means
  \fill[targetgreen!5,rounded corners=4pt,draw=targetgreen!25] (-0.8, -1.0) rectangle (12.2, 1.4);
  \node[lanehead,text=targetgreen!90!black] at (-0.6, 1.0) {Lane 3: Parallel Registers $Z_r \gets Z_r + \Delta_t h_r(e_t)$ and Median of Means};

  \node[reg] (z1) at (1.5, 0.3) {$Z_1$};
  \node[reg] (z2) at (5.5, 0.3) {$Z_2$};
  \node[reg] (zd) at (9.5, 0.3) {$Z_d$};

  \draw[->,quotientorange!80!black,thick] (h1) -- (z1);
  \draw[->,quotientorange!80!black,thick] (h2) -- (z2);
  \draw[->,quotientorange!80!black,thick] (hd) -- (zd);

  \node[font=\scriptsize,text=gray!80!black] at (1.5, -0.4) {$Y_1 = Z_1^2$};
  \node[font=\scriptsize,text=gray!80!black] at (5.5, -0.4) {$Y_2 = Z_2^2$};
  \node[font=\scriptsize,text=gray!80!black] at (9.5, -0.4) {$Y_d = Z_d^2$};

  \draw[->,targetgreen!80!black,very thick] (5.5, -0.6) -- (5.5, -0.9)
    node[below,font=\footnotesize\bfseries,text=targetgreen!90!black]
    {$\widehat{F}_2 = \operatorname{median}(\text{mean}(Y_{r,1..w})) \approx (1\pm\eps) F_2$};
\end{tikzpicture}
\caption{The Alon--Matias--Szegedy (AMS) sketch architecture for graph frequency moments. Lane~1: stream tokens arrive with edge identifier and sign ($\Delta_t = \pm 1$). Lane~2: $d$ independent $4$-wise hash functions assign uniform signs $h_r(e) \in \{-1, +1\}$ without storing lookup tables. Lane~3: scalar accumulators $Z_r$ track the random projection. The squared values $Y_r = Z_r^2$ are averaged within rows of width $w = \Oh(1/\eps^2)$ to suppress variance, and the median across $d = \Oh(\log(1/\delta))$ rows yields a $(1\pm\eps)$-approximation with probability at least $1-\delta$.}
\label{ch40:fig-ams-sketch}
\end{figure}

\section{What small space cannot do}
\label{ch40:lower}

The positive results above are tightly bounded by impossibility theorems.
The standard tool is \emph{communication complexity}: split the stream between
two players, Alice (the prefix) and Bob (the suffix); a streaming algorithm
using $s$ bits induces a protocol where Alice sends Bob her $s$-bit memory
state, so any $s$-bit one-pass algorithm yields an $s$-bit one-way protocol.
Lower-bounding the communication then lower-bounds $s$.

\begin{theorem}[Connectivity needs $\Om(n)$ space]\label{ch40:thm-lb-conn}
Any (even randomized, constant-error) one-pass streaming algorithm that
decides connectivity uses $\Om(n)$ bits.
\end{theorem}

\begin{proof}[Proof idea]
Reduce from one-way \textsc{Index}. Alice holds a string $y\in\set{0,1}^N$
with $N=\Th(n)$, Bob holds an index $k$, and Bob must output $y_k$;
one-way communication forces $\Om(N)$ bits. Encode $y$ as the presence/absence
of $N$ candidate edges incident to a gadget so that the final graph is
connected iff $y_k=1$. Alice streams her edges and passes the memory state;
Bob streams the query edges; the connectivity answer reveals $y_k$. Since the
protocol transmits exactly the memory state, $s=\Om(N)=\Om(n)$.
\end{proof}

\begin{figure}[H]
\centering
\begin{tikzpicture}[font=\small,
  vtx/.style={circle,draw=nodeblue,fill=blue!6,minimum size=6mm,inner sep=0pt}]
\node[draw=sourceblue,thick,rounded corners,minimum width=3.2cm,
  minimum height=2.6cm,label=above:{\scriptsize Alice (prefix of $\sigma$)}]
  (A) at (0,0) {};
\node[draw=targetgreen,thick,rounded corners,minimum width=2.6cm,
  minimum height=2.6cm,label=above:{\scriptsize Bob (suffix of $\sigma$)}]
  (B) at (5.4,0) {};
\draw[-{Stealth[length=2.5mm]},thick]
  ($(A.east)+(0,0.2)$) -- node[above,font=\scriptsize]{$s$-bit state}
  ($(B.west)+(0,0.2)$);
\node[draw=none,fill=none,font=\scriptsize] at (2.7,-1.0)
  {connected? $\;\Leftrightarrow\; y_k=1$};
\end{tikzpicture}
\caption{The communication reduction behind
Theorem~\ref{ch40:thm-lb-conn}. Alice encodes her bitstring $y$ as the prefix
of the edge stream and forwards only the algorithm's $s$-bit memory state; Bob
appends query edges for index $k$ and reads $y_k$ off the connectivity answer.
A protocol this cheap forces $s=\Om(N)=\Om(n)$.}
\label{ch40:fig-lb}
\end{figure}

So semi-streaming connectivity ($\Oh(n\log n)$) is optimal up to the
logarithmic factor. Distances are dramatically worse:

\begin{theorem}[No exact distances, and stretch needs space]\label{ch40:thm-lb-dist}
Any one-pass algorithm that, for a designated pair $(s,t)$, distinguishes
$\dist(s,t)\le d$ from $\dist(s,t)>d$ exactly requires $\Om(n^2)$ bits in the
worst case --- essentially the whole graph. More generally, a one-pass
$(2k-1)$-spanner-quality stretch bound forces $\Om(n^{1+1/k})$ space, matching
Algorithm~\ref{ch40:alg-spanner} up to logs.
\end{theorem}

\begin{intuitionbox}
The spanner lower bound is the Moore bound (Lemma~\ref{ch40:lem-moore}) read
backwards. There exist graphs of girth $> 2k$ with $\Th(n^{1+1/k})$ edges in
which \emph{every} edge is essential: delete any one and some pair's distance
jumps. To certify stretch $< 2k+1$ you must remember $\Om(n^{1+1/k})$ bits,
because forgetting any edge of such a graph is detectable. Stretch and space
trade off along exactly the same curve the greedy spanner achieves --- the
algorithm is optimal, not merely good.
\end{intuitionbox}

\begin{theorem}[Triangle detection needs $\Om(n^2)$ space]\label{ch40:thm-lb-tri}
Distinguishing $T=0$ from $T\ge1$ in one pass requires $\Om(n^2)$ bits.
Consequently any multiplicative $(1\pm\eps)$ approximation of $T$ needs space
growing as $T$ shrinks; no algorithm in $o(n^2)$ space can certify
triangle-freeness.
\end{theorem}

\begin{proof}[Proof idea]
Reduce from two-party \textsc{Disjointness}: Alice's set and Bob's set
intersect iff a specific triangle exists in a bipartite-plus-query gadget.
\textsc{Disjointness} needs $\Om(N)=\Om(n^2)$ communication even for
randomized protocols, and again the protocol cost equals the memory state, so
$s=\Om(n^2)$.
\end{proof}

\begin{trapbox}
A multiplicative approximation of $T$ \emph{cannot} tell $T=0$ from $T=1$ ---
$(1\pm\eps)\cdot0=0$, yet you would need to detect the single triangle, which
Theorem~\ref{ch40:thm-lb-tri} forbids in small space. This is why
Algorithm~\ref{ch40:alg-tri}'s sample complexity blows up as $T\to0$: the
estimator is not broken, the problem is provably hard there. Always read an
approximation guarantee together with what it silently excludes.
\end{trapbox}

\subsection{The tradeoffs, laid bare}

There is no free lunch; there are only well-mapped exchanges. Pin them down:

\begin{center}
\begin{tabular}{lllll}
\toprule
problem & model & space & passes & answer quality \\
\midrule
connectivity & insertion & $\Oh(n\log n)$ & 1 & exact \\
connectivity & turnstile & $\Oh(n\log^3 n)$ & 1 & exact whp \\
bipartiteness & insertion & $\Oh(n\log n)$ & 1 & exact \\
$(2k{-}1)$-spanner & insertion & $\Oh(n^{1+1/k}\log n)$ & 1 & stretch $2k{-}1$ \\
exact distance & either & $\Om(n^2)$ & 1 & exact (impossible small) \\
triangle count & insertion & $\widetilde\Oh\!\big(\tfrac{m\Delta}{\eps^2 T}\big)$ & 1 & $(1{\pm}\eps)$ whp \\
triangle $T\!>\!0$? & either & $\Om(n^2)$ & 1 & exact (impossible small) \\
\bottomrule
\end{tabular}
\end{center}

Three axes are in tension, and you can only bend one by spending another:

\begin{itemize}
\item \textbf{Space vs.\ accuracy.} Spanners trade more space ($n^{1+1/k}$) for
less stretch ($2k-1$); triangle estimators trade more samples for smaller
$\eps$. The frontier curves are matched by lower bounds --- you cannot cheat
them.
\item \textbf{Space vs.\ passes.} Many problems collapse to $\Oh(n)$ space if
you allow $\Oh(\log n)$ passes (e.g.\ exact shortest paths via repeated
Bellman--Ford-style sweeps). One pass is the strict regime; each extra pass is
a real-world cost (re-reading petabytes) but a real algorithmic lever.
\item \textbf{Insertion vs.\ turnstile.} Deletions force linearity, which costs
extra log factors and forces randomization (the sketches fail with small
probability). Insertion-only lets you be greedy and deterministic. If your
application never deletes, do not pay the sketching tax.
\end{itemize}

\begin{aestheticbox}
The recurring shape of every result here is the same: \emph{a graph's useful
structure is near-linear-size and linearly compressible}, while its
\emph{exact} fine structure (precise distances, the presence of one specific
triangle) is irreducibly quadratic. Semi-streaming is the science of living on
the right side of that divide --- keeping one number per vertex, exploiting cut
and boundary algebra, and accepting approximation exactly where the lower
bounds say you must. It is the same compressibility that powers spectral
sparsifiers and Laplacian solvers, viewed through the keyhole of a single
left-to-right scan.
\end{aestheticbox}

\section{Historical Context, Controversies, and the Road Ahead}

\begin{historybox}
The mathematical foundation of streaming algorithms was established in $1996$ by Noga Alon, Yossi Matias, and Mario Szegedy in their Gödel-prize-winning paper on frequency moments, showing that statistical properties of massive streams could be tracked using tiny randomized sketches. In the early $2000$s, as web graphs and internet traffic exploded, computer scientists realized that classical graph theory's core assumption---random access to an entire adjacency matrix in RAM---had broken down. In $2005$, Joan Feigenbaum, Sampath Kannan, Andrew McGregor, Siddharth Suri, and Jian Zhang introduced the \emph{semi-streaming model}, proposing an $O(n \operatorname{polylog} n)$ space budget: storing a few numbers per vertex is feasible on web graphs while storing edges is impossible.

The decisive breakthrough for dynamic streams arrived in $2012$, when Kook Jin Ahn, Sudipto Guha, and Andrew McGregor (AGM) introduced \emph{graph sketching}. AGM demonstrated that by representing a graph through its signed vertex-edge incidence matrix and computing linear $L_0$ sketches over its rows, one could maintain spanning forests, connectivity, and cuts in a single pass under arbitrary edge insertions and deletions.

The transition to streaming exposed the limits of standard greedy approaches. In an insertion stream, keeping a spanning forest greedily by discarding cycle-forming edges works perfectly; but in a turnstile stream where edges can be deleted, greedy algorithms fail catastrophically because discarded edges cannot be retrieved when a tree edge is removed. Overcoming this required linear sketches that store numerical aggregates that automatically undo themselves when deleted. Furthermore, two-party communication complexity proved that computing exact shortest paths in a single pass requires $\Omega(n^2)$ bits: exact distance is an irreducibly quadratic property, forcing algorithm designers to embrace metric spanners with stretch $2k-1$.
\end{historybox}

\subsection{What We Can Now Do vs.\ Honest Downsides and the Road Ahead}
\begin{connectionbox}
\textbf{Capabilities Unlocked}:
\begin{itemize}
  \item We can process massive graphs with billions of edges in a single sequential pass using working memory proportional only to vertex count ($O(n \operatorname{polylog} n)$).
  \item Linear sketching (AGM) allows maintaining spanning forests and connectivity under arbitrary edge insertions and deletions (turnstile streams).
  \item Metric spanners provide provable $(2k-1)$ distance approximations with $O(n^{1+1/k})$ edges in an insertion stream.
\end{itemize}

\textbf{Honest Downsides and Limitations}:
\begin{itemize}
  \item \emph{Inability to solve fine-grained problems}: Exact distances, exact triangle counts, and exact subgraph detection provably require $\Omega(n^2)$ space in one pass; exact answers cannot be obtained in sublinear memory.
  \item \emph{Randomization failure risk}: AGM sketches succeed with high probability (e.g., $1 - 1/n^c$). In rare adversarial configurations, hash collisions can cause false negative connectivity reports.
  \item \emph{Per-token computational overhead}: While insertion-only streams process tokens in nanoseconds via union-find, turnstile linear sketching updates $O(\log^3 n)$ hash registers for every edge event, reducing throughput.
\end{itemize}

\textbf{Where We Go Next}:
\begin{itemize}
  \item In Chapter~41 (Laplacian Solvers and Electrical Flows), we see how linear sketching and edge sampling concepts culminate in continuous numerical graph algorithms. Spielman and Srivastava use effective resistance sampling to construct $(1 \pm \eps)$ \emph{spectral sparsifiers} with only $O(n \log n / \eps^2)$ edges, unlocking near-linear time Laplacian solvers.
  \item In Chapter~42 (Appendix: Math and Data Structures), we review the foundational data structures (DSU, segment trees) and linear algebra (Courant--Fischer min-max, matrix concentration) supporting these algorithms.
\end{itemize}
\end{connectionbox}

\section{Exercises}
\label{ch40:exercises}

\begin{exercise}[Implement: components and a witness]
Implement Algorithm~\ref{ch40:alg-forest} with union by rank and path
compression. Extend it so that, on a disconnected input, it outputs an actual
vertex pair $(u,v)$ in different components as a certificate of
disconnection. Argue your extension stays within $\Oh(n\log n)$ bits.
\end{exercise}

\begin{exercise}[Prove: bipartiteness parity bits]
Prove that the parity-bit augmentation in
Section~\ref{ch40:connectivity} correctly reports non-bipartiteness exactly
when an odd cycle exists. Show that maintaining $\phi(v)$ under union--find
merges (flipping a whole subtree's parities when roots merge with conflicting
orientation) is correct, and bound the total parity-update work.
\end{exercise}

\begin{exercise}[Break: greedy forest under deletions]
Exhibit a turnstile stream on $4$ vertices on which
Algorithm~\ref{ch40:alg-forest}, naively re-run with deletions removing edges
from $F$, reports the wrong number of components. Explain precisely which step
of the correctness proof in Section~\ref{ch40:connectivity} fails, and why
linearity in Section~\ref{ch40:sketching} repairs it.
\end{exercise}

\begin{exercise}[Spanner stretch is tight]
For $t=3$ ($k=2$), construct a graph and a stream order on which
Algorithm~\ref{ch40:alg-spanner} produces an $H$ where some original edge
$\set{u,v}$ has $\dist_H(u,v)=3$ exactly. Conclude the stretch bound $2k-1$ is
achieved, not merely upper-bounded.
\end{exercise}

\begin{exercise}[Moore bound, the details]
Fill in the cleaning step of Lemma~\ref{ch40:lem-moore}: prove that any graph
with average degree $d$ contains a subgraph with minimum degree $\ge d/2$
(repeatedly delete vertices of degree $< d/2$ and bound how many edges leave).
Then complete the BFS-ball counting to get $\abs{E}=\Oh(n^{1+1/k})$.
\end{exercise}

\begin{exercise}[$1$-sparse recovery]
A vector $x\in\R^N$ has at most one nonzero coordinate. You are told only
$P=\sum_i x_i$ and $Q=\sum_i i\,x_i$. Show how to recover the index and value
of the nonzero coordinate, and explain what goes wrong if the support has size
two. Then describe how a random fingerprint slot detects (whp) a support of
size $\ge 2$ so the sampler can reject.
\end{exercise}

\begin{exercise}[Why fresh sketches per round]
In Algorithm~\ref{ch40:alg-agm}, suppose you reuse a single $L_0$ sketch
across all $R$ Bor\r{u}vka rounds. Construct (or carefully argue) a scenario
where reuse causes a correlated failure: the round-$r$ query set
$\set{S\in\mathcal{C}}$ depends on the sketch's own randomness from earlier
rounds, breaking the uniformity guarantee. How many independent sketches are
truly necessary, and why $\Th(\log n)$?
\end{exercise}

\begin{exercise}[Triangle estimator variance]
For Algorithm~\ref{ch40:alg-tri}, compute $\mathrm{Var}[\widehat T]$ for a
single copy in terms of $m,n,T$, and the maximum number $\Delta$ of triangles
sharing an edge. Deduce the number of parallel copies needed for a
$(1\pm\eps)$ estimate with probability $1-\delta$, and verify it matches the
$\widetilde\Oh(m\Delta/(\eps^2 T))$ figure quoted in
Section~\ref{ch40:triangles}.
\end{exercise}

\begin{exercise}[Lower bound from \textsc{Index}]
Write out the reduction behind Theorem~\ref{ch40:thm-lb-conn} in full: define
the edge gadget encoding $y\in\set{0,1}^N$, specify Bob's query edges for
index $k$, and prove the final graph is connected iff $y_k=1$. State exactly
where the one-way communication lower bound of $\Om(N)$ enters.
\end{exercise}

\begin{exercise}[Passes vs.\ space]
Show that with $\Oh(\log n)$ passes and $\Oh(n)$ space you can compute exact
single-source BFS distances in an unweighted streamed graph (one BFS layer per
pass). Contrast with the one-pass $\Om(n^2)$ bound of
Theorem~\ref{ch40:thm-lb-dist} and explain, in your own words, what extra
passes buy.
\end{exercise}

\begin{exercise}[Extend: connectivity sketch to MST weight]
Sketch (informally) how to extend the AGM connectivity sketch of
Algorithm~\ref{ch40:alg-agm} to estimate the weight of a minimum spanning
tree of a turnstile-streamed weighted graph, using the
``count components at each weight threshold'' telescoping identity
$w(\mathrm{MST}) = \sum_{i} \bigl(\,\text{\#components using only edges of
weight} \le i\,\bigr) - (\,\dots\,)$. Which pieces of this chapter does it
reuse?
\end{exercise}
